\section{내용, 원시다항식과 가우스 보조정리}
\label{sec:gauss-lemma}

\begin{learninggoal}
정수계수 다항식의 공통인수를 분리하고, $\Z[X]$에서의 인수분해와 $\Q[X]$에서의 인수분해가 어떻게 연결되는지 이해한다.
\end{learninggoal}

이 절에서는 $R$을 UFD, $K=Q(R)$를 그 분수체라 한다.

\begin{definition}
영이 아닌 다항식
\[
  f=a_0+a_1X+\cdots+a_nX^n\in R[X]
\]
의 계수들의 최대공약수를 $f$의 \emph{내용}이라 하고 $\operatorname{cont}(f)$로 쓴다. 단위원배의 차이는 허용한다. 내용이 단위원이면 $f$를 \emph{원시다항식}이라 한다.
\end{definition}

예를 들어 $6X^3-9X+3$의 내용은 $3$이고
\[
  6X^3-9X+3=3(2X^3-3X+1)
\]
에서 괄호 안의 다항식은 원시다항식이다.

\begin{lemma}[가우스 보조정리]
원시다항식 두 개의 곱은 원시다항식이다. 더 일반적으로 영이 아닌 $f,g\in R[X]$에 대해
\[
  \operatorname{cont}(fg)\sim
  \operatorname{cont}(f)\operatorname{cont}(g),
\]
여기서 $\sim$는 단위원배를 제외하고 같다는 뜻이다.
\end{lemma}

\begin{proofidea}
곱의 모든 계수가 어떤 소원소 $p$로 나누어진다고 가정한다. 계수를 $p$로 줄이면 두 원시다항식은 모두 영이 아닌 다항식이지만 그 곱은 영다항식이 된다. 그러나 $(R/(p))[X]$는 정역이므로 모순이다.
\end{proofidea}

\begin{proof}
원시인 $f,g$의 곱이 원시가 아니라고 하자. UFD에서 어떤 소원소 $p$가 $fg$의 모든 계수를 나눈다. 계수를 $(p)$로 줄이는 준동형
\[
  R[X]\to(R/(p))[X]
\]
을 생각하면, 원시성 때문에 $f,g$의 상은 각각 영이 아니다. $p$가 소원소이므로 $(p)$는 소아이디얼이고 $R/(p)$는 정역이다. 따라서 $(R/(p))[X]$도 정역인데 $fg$의 상은 영이므로 모순이다.

일반적인 $f,g$를
\[
  f=\operatorname{cont}(f)f_0,
  \qquad
  g=\operatorname{cont}(g)g_0
\]
로 쓰자. 여기서 $f_0,g_0$는 원시다항식이다. 방금 증명한 결과로 $f_0g_0$도 원시이므로 $fg$의 내용은 $\operatorname{cont}(f)\operatorname{cont}(g)$와 동반이다.
\end{proof}

\begin{lemma}[원시다항식의 유리수배]
$P\in R[X]$가 원시이고 $c\in K^\times$라 하자. $cP\in R[X]$이면 $c\in R$이다. 더 나아가 $cP$도 원시이면 $c\in R^\times$이다.
\end{lemma}

\begin{proof}
$c=a/b$로 쓰되 $a,b\in R$, $b\ne0$이고 $a,b$가 서로소라고 하자. $cP$의 모든 계수가 $R$에 속하므로 $b$는 $a$와 $P$의 각 계수의 곱을 나눈다. 만약 어떤 소원소 $p$가 $b$를 나누면 $p\nmid a$이고, 소원소의 성질에 의해 $p$가 $P$의 모든 계수를 나누게 된다. 이는 $P$의 원시성에 모순이다. 따라서 $b$는 단위원이고 $c\in R$이다.

또한 $cP$가 원시인데 $c$가 비단위원이면, $c$의 어떤 소인수가 $cP$의 모든 계수를 나누어 다시 모순이다. 따라서 $c$는 단위원이다.
\end{proof}

\begin{theorem}[가우스의 기약성 정리]
양의 차수인 원시다항식 $f\in R[X]$에 대해 다음은 동치이다.
\begin{enumerate}[label=(\roman*)]
  \item $f$는 $R[X]$에서 기약이다.
  \item $f$는 $K[X]$에서 기약이다.
\end{enumerate}
\end{theorem}

\begin{proofidea}
$K[X]$에서의 인수분해가 있으면 각 인수의 분모를 없애고 내용을 제거하여 $R[X]$의 원시 인수들로 바꾼다. 앞의 보조정리는 남는 상수가 반드시 $R$의 단위원임을 보장한다.
\end{proofidea}

\begin{proof}
먼저 $f$가 $R[X]$에서 가약이라고 하자. 원시성 때문에 비자명한 인수분해 $f=gh$에서 어느 인수도 비단위원 상수일 수 없다. 실제로 그런 상수의 소인수는 $f$의 모든 계수를 나누게 된다. 따라서 $g,h$는 모두 양의 차수이고, 같은 식은 $K[X]$에서도 비자명한 인수분해이다. 그러므로 (ii)는 (i)를 함의한다.

반대로 $f=gh$가 $K[X]$에서의 비자명한 인수분해라고 하자. 적당한 $u,v\in K^\times$를 택하여
\[
  g_0=ug,
  \qquad
  h_0=vh
\]
를 $R[X]$의 원시다항식으로 만들 수 있다. 그러면 어떤 $c\in K^\times$에 대해
\[
  f=cg_0h_0
\]
이다. 가우스 보조정리에 의해 $g_0h_0$는 원시이고, $f$도 원시이다. 앞의 보조정리를 적용하면 $c\in R^\times$이다. 따라서
\[
  f=(cg_0)h_0
\]
는 $R[X]$에서 양의 차수인 두 다항식의 곱으로 분해된다. 그러므로 (i)는 (ii)를 함의한다.
\end{proof}

\begin{corollary}[가우스 정리]
$R$이 UFD이면 $R[X]$도 UFD이다. 귀납적으로
\[
  R[X_1,\dots,X_n]
\]
도 UFD이다.
\end{corollary}

\begin{proof}
영이 아닌 비단위원 $f\in R[X]$를
\[
  f=c f_0
\]
로 쓰자. 여기서 $c=\operatorname{cont}(f)\in R$이고 $f_0$는 원시이다. $c$는 $R$에서 소원소들의 곱으로 분해된다. 한편 $K[X]$는 유클리드 정역이므로 UFD이고, $f_0$는 $K[X]$에서 기약다항식들의 곱으로 분해된다. 각 인수에 적절한 상수를 곱하여 $R[X]$의 원시다항식으로 만들고, 가우스의 기약성 정리를 적용하면 이 인수들은 $R[X]$에서도 기약이다. 원시다항식의 곱이 원시라는 사실과 위의 유리수배 보조정리를 사용하면 남는 전체 상수는 $R$의 단위원에 흡수된다. 따라서 $f$는 $R[X]$의 기약원소들의 곱으로 분해된다.

$R[X]$의 상수 기약원소는 $R$의 기약원소이고, UFD인 $R$에서는 소원소이다. 한편 $R[X]$의 양의 차수 기약원소는 반드시 원시이다. 내용이 비단위원이면 내용과 원시 부분의 곱이 비자명한 인수분해를 주기 때문이다. 따라서 기약원소가 소원소임을 보이기 위해 다음 두 경우만 확인하면 충분하다.

먼저 $R$의 소원소 $p$에 대해
\[
  R[X]/(p)\cong(R/(p))[X]
\]
는 정역이므로 $p$는 $R[X]$에서도 소원소이다. 다음으로 원시 기약다항식 $q\in R[X]$는 $K[X]$에서 기약이고 따라서 $K[X]$에서 소원소이다. $q\mid FG$가 $R[X]$에서 성립하면 $K[X]$에서 $q\mid F$ 또는 $q\mid G$이다. 예를 들어 $F=qh$인 $h\in K[X]$가 존재하면 $h=c h_0$으로 원시 부분을 분리할 수 있고, $F=c(qh_0)\in R[X]$와 $qh_0$의 원시성에서 유리수배 보조정리에 의해 $c\in R$을 얻는다. 따라서 $h\in R[X]$이고 $q\mid F$가 $R[X]$에서도 성립한다. 즉 모든 기약원소가 소원소이므로 $R[X]$는 UFD이다.

여러 변수의 경우
\[
  R[X_1,\dots,X_n]
  =R[X_1,\dots,X_{n-1}][X_n]
\]
에 같은 결과를 반복 적용한다.
\end{proof}

\begin{example}
$f=2X^3+3X+1\in\Z[X]$는 원시다항식이다. 따라서 $f$가 $\Q[X]$에서 기약인지 조사할 때 정수계수 인수분해만 고려해도 충분하다. 분수계수 인수가 갑자기 새롭게 나타나는 일은 없다.
\end{example}

\begin{pitfall}
가우스 정리는 ``$\Z[X]$에서 기약이면 모든 더 큰 체에서도 기약''이라는 뜻이 아니다. 예를 들어 $X^2+1$은 $\Q[X]$에서 기약이지만 $\C[X]$에서는 분해된다. 정리는 UFD $R$과 바로 그 분수체 $Q(R)$ 사이를 비교한다.
\end{pitfall}

\begin{keyidea}
내용은 계수에 숨어 있는 상수 인수이고, 원시 부분은 다항식 자체의 인수분해 구조를 담는다. 가우스 보조정리는 이 두 층이 서로 간섭하지 않음을 말한다.
\end{keyidea}
